% ==================================================================
% id:       func-006
% topic:    פונקציות
% subtopic: תחום הגדרה ותכונות של פונקציות
% solution: full
% ==================================================================
\textbf{לגבי כל אחת מהפונקציות הבאות, קבעו האם הפונקציה זוגית, והאם היא אי-זוגית.}
\begin{enumerate}\bMs{2}
\item $f(x) = \frac{x}{x^2 -1}$
\begin{solbox}
אי-זוגית. תחום ההגדרה הוא כל $x\neq\pm1$, והוא סימטרי ביחס ל-$0$. לכל $x$ בתחום מתקיים
$$.f(-x)=\frac{-x}{(-x)^2-1}=-\frac{x}{x^2-1}=-f(x)$$
הפונקציה אינה זוגית, כי למשל
$f(2)=\frac23$
ואילו
$f(-2)=-\frac23$.
\end{solbox}
\item $f(x) = x^4 - 3x^2 + 1$
\begin{solbox}
זוגית. הפונקציה מוגדרת לכל $x$ ממשי, ומתקיים
$$.f(-x)=(-x)^4-3(-x)^2+1=x^4-3x^2+1=f(x)$$
הפונקציה אינה אי-זוגית, כי למשל
$f(1)=-1$
וגם
$f(-1)=-1$,
ולכן
$f(-1)\neq -f(1)$.
\end{solbox}
\item $f(x) = \sin (2x)$
\begin{solbox}
אי-זוגית. הפונקציה מוגדרת לכל $x$ ממשי, ומכיוון ש-$\sin$ אי-זוגית מתקיים
$$.f(-x)=\sin(-2x)=-\sin(2x)=-f(x)$$
הפונקציה אינה זוגית, כי למשל
$f\lr{\frac\pi4}=1$
ואילו
$f\lr{-\frac\pi4}=-1$.
\end{solbox}
\item $f(x) = \tan (x + 1)$
\hint{כדי להראות שתכונה לא מתקיימת מספיק למצוא נקודה אחת שבה היא נכשלת. השוו בין הערכים של הפונקציה ב-$x=1$ וב-$x=-1$.}
\begin{solbox}
לא זוגית ולא אי-זוגית. הנקודות $1$ ו-$-1$ שתיהן בתחום ההגדרה, ומתקיים
$$.f(-1)=\tan 0=0,\qquad f(1)=\tan 2\neq 0$$
(אכן
$\tan 2\neq0$,
כי
$\tan$
מתאפסת רק בכפולות שלמות של $\pi$.)
לכן
$f(-1)\neq f(1)$,
והפונקציה אינה זוגית, וגם
$f(-1)\neq -f(1)$,
והפונקציה אינה אי-זוגית.
\end{solbox}
\item $f(x)=0$
\begin{solbox}
גם זוגית וגם אי זוגית, שכן לכל $x$ ממשי מתקיים
$f(-x) = 0 = f(x) = -f(x)$
\end{solbox}
\item  $\ds{f(x)=\frac{x^2}{\sqrt{x^4+1}}}$
\begin{solbox}
ראשית נשים לב שהפונקציה מוגדרת לכל הישר הממשי. והיא זוגית כמנה של פונקציות זוגיות.
\end{solbox}
\item  $f(x)=x+1$
\begin{solbox}
לא זוגית ולא אי זוגית דוגמה נגדית:
$f(1)=2 , f(-1)=0$
\end{solbox}
\item  $f(x)=x^3$
\begin{solbox}
אי זוגית,
מאחר וחזקה אי זוגית היא פונקציה אי זוגית.
\end{solbox}
\item  $\ds{f(x)=\log_2\lr{\frac{1-x^3}{1+x^3}}}$
\hint{חשבו את $f(-x)$ והיעזרו בתכונות הלוגריתם.}
\begin{solbox}
הפונקציה מוגדרת כאשר
$\frac{1-x^3}{1+x^3}>0$,
כלומר עבור
$-1<x<1$,
וזהו תחום סימטרי סביב $0$.
לכל $x$ בתחום מתקיים
$$f(-x)=\log_2\lr{\frac{1+x^3}{1-x^3}}=\log_2\lr{\lr{\frac{1-x^3}{1+x^3}}^{-1}}=-\log_2\lr{\frac{1-x^3}{1+x^3}}=-f(x),$$
ולכן הפונקציה אי-זוגית.
\end{solbox}
\item $f(x)=\sin^2(x)+\cos(x)$
\hint{חשבו את $f(-x)$ והיעזרו בזוגיות או באי-הזוגיות של $\sin$ ושל $\cos$.}
\begin{solbox}
פונקציה זוגית,
כסכום של פונקציות זוגיות.
ידוע כי
$\cos(x)$
זוגית
ו
$sin(x)$
אי זוגית ולכן
$sin^2(x)$
זוגית
כי
\\
$sin^2(-x)=(sin(-x))^2=(-sin(x))^2=sin^2(x)$
\end{solbox}
\eMs
\end{enumerate}
